% GENERATED FILE. Edit canonical lessons, then run npm run generate:books.
% canonical-source-hash: fnv1a64:0651974941fbbad8
% canonical-lessons: FR-C419-decoller, FR-C419-decollage, FR-C419-atterrissage, FR-C419-atterrir, FR-C419-survoler

\chapter{Décollage et atterrissage}
\label{ch:fr-decollage-et-atterrissage}
\hlchaptermodality{419}

\begin{chapteropening}
\textbf{By the end of this chapter:} \emph{I can talk about take-off and landing, and flying over places.}

Complete the last lesson of chapter 419: I can talk about take-off and landing, and flying over places.
\end{chapteropening}

\section[décoller]{décoller — to take off; to unstick}

\label{lesson:FR-C419-decoller}



\begin{quote}
Before the new one: say the French for the porthole, then the French for the aisle.
\end{quote}

\subsection*{You'll want to know: décoller}
\textbf{décoller} — \textquotedblleft{}to take off; to unstick\textquotedblright{}.

\begin{quote}
L'avion a décollé avec une heure de retard.
\end{quote}

\begin{quote}
\emph{The plane took off an hour late.}
\end{quote}

Also \emph{décoller une étiquette} — to peel off a label.

\begin{cousinweb}
\emph{décoller} is built from \emph{dé-} 'un-' and \emph{coller} 'to stick', from \emph{colle} 'glue', which goes back to Greek \emph{kólla} (EtymDB): the plane comes unstuck from the ground.
\end{cousinweb}

\begin{grammarlens}[title={What you've built}]
One more word for telling, explaining and describing at B1.
\end{grammarlens}

\subsection*{Your turn}
\begin{itemize}
\raggedright
  \item \emph{Say it:} \emph{décoller}
  \item \emph{Say it:} \emph{décoller}, once more
  \item \emph{Say it:} a sentence of your own with \emph{décoller}
  \item \emph{Recall:} say the French for the porthole, then the French for the aisle, then say \emph{décoller} again
\end{itemize}

\begin{tcolorbox}[breakable,colback=teal!4,colframe=teal!35!black,title={Before you move on}]
What does décoller mean? (\textquotedblleft{}To take off; to unstick\textquotedblright{}.) Say it once more.
\end{tcolorbox}
\section[le décollage]{le décollage — the take-off}

\label{lesson:FR-C419-decollage}



\begin{quote}
Before the new one: say the French for the aisle, then the French for to take off.
\end{quote}

\subsection*{You'll want to know: le décollage}
\textbf{le décollage} — \textquotedblleft{}the take-off\textquotedblright{}.

\begin{quote}
Au décollage, les enfants regardent par la fenêtre.
\end{quote}

\begin{quote}
\emph{At take-off the children look out of the window.}
\end{quote}

Built on \emph{décoller} with \emph{-age}.

\begin{cousinweb}
\emph{décollage} is built on \emph{décoller} with \emph{-age}. \emph{Décoller} is \emph{dé-} 'un-' plus \emph{coller} 'to stick', from \emph{colle} 'glue', from Greek \emph{kólla} (EtymDB).
\end{cousinweb}

\begin{grammarlens}[title={What you've built}]
One more word for telling, explaining and describing at B1.
\end{grammarlens}

\subsection*{Your turn}
\begin{itemize}
\raggedright
  \item \emph{Say it:} \emph{le décollage}
  \item \emph{Say it:} \emph{le décollage}, once more
  \item \emph{Say it:} a sentence of your own with \emph{le décollage}
  \item \emph{Recall:} say the French for the aisle, then the French for to take off, then say \emph{le décollage} again
\end{itemize}

\begin{tcolorbox}[breakable,colback=teal!4,colframe=teal!35!black,title={Before you move on}]
What does le décollage mean? (\textquotedblleft{}The take-off\textquotedblright{}.) Say it once more.
\end{tcolorbox}
\section[l'atterrissage]{l'atterrissage — the landing (of a plane)}

\label{lesson:FR-C419-atterrissage}



\begin{quote}
Before the new one: say the French for to take off, then the French for the take-off.
\end{quote}

\subsection*{You'll want to know: l'atterrissage}
\textbf{l'atterrissage} — \textquotedblleft{}the landing (of a plane)\textquotedblright{}.

\begin{quote}
L'atterrissage a été un peu brusque.
\end{quote}

\begin{quote}
\emph{The landing was a bit rough.}
\end{quote}

A masculine noun.

\begin{cousinweb}
\emph{atterrissage} is built on \emph{atterrir} 'to land' with \emph{-age}; \emph{atterrir} comes from \emph{terre} 'land', from Latin \emph{terra} (Lewis \& Short). EtymDB does not list the noun.
\end{cousinweb}

\begin{grammarlens}[title={What you've built}]
One more word for telling, explaining and describing at B1.
\end{grammarlens}

\subsection*{Your turn}
\begin{itemize}
\raggedright
  \item \emph{Say it:} \emph{l'atterrissage}
  \item \emph{Say it:} \emph{l'atterrissage}, once more
  \item \emph{Say it:} a sentence of your own with \emph{l'atterrissage}
  \item \emph{Recall:} say the French for to take off, then the French for the take-off, then say \emph{l'atterrissage} again
\end{itemize}

\begin{tcolorbox}[breakable,colback=teal!4,colframe=teal!35!black,title={Before you move on}]
What does l'atterrissage mean? (\textquotedblleft{}The landing (of a plane)\textquotedblright{}.) Say it once more.
\end{tcolorbox}
\section[atterrir]{atterrir — to land}

\label{lesson:FR-C419-atterrir}



\begin{quote}
Before the new one: say the French for the take-off, then the French for the landing.
\end{quote}

\subsection*{You'll want to know: atterrir}
\textbf{atterrir} — \textquotedblleft{}to land\textquotedblright{}.

\begin{quote}
Nous avons atterri sous la pluie.
\end{quote}

\begin{quote}
\emph{We landed in the rain.}
\end{quote}

Conjugated like \emph{finir}: \emph{nous atterrissons}.

\begin{cousinweb}
\emph{atterrir} is built from \emph{à} and \emph{terre} 'land', which is inherited from Latin \emph{terra} (Lewis \& Short); EtymDB lists the verb without a source.
\end{cousinweb}

\begin{grammarlens}[title={What you've built}]
One more word for telling, explaining and describing at B1.
\end{grammarlens}

\subsection*{Your turn}
\begin{itemize}
\raggedright
  \item \emph{Say it:} \emph{atterrir}
  \item \emph{Say it:} \emph{atterrir}, once more
  \item \emph{Say it:} a sentence of your own with \emph{atterrir}
  \item \emph{Recall:} say the French for the take-off, then the French for the landing, then say \emph{atterrir} again
\end{itemize}

\begin{tcolorbox}[breakable,colback=teal!4,colframe=teal!35!black,title={Before you move on}]
What does atterrir mean? (\textquotedblleft{}To land\textquotedblright{}.) Say it once more.
\end{tcolorbox}
\section[survoler]{survoler — to fly over}

\label{lesson:FR-C419-survoler}



\begin{quote}
Before the new one: say the French for the landing, then the French for to land.
\end{quote}

\subsection*{You'll want to know: survoler}
\textbf{survoler} — \textquotedblleft{}to fly over\textquotedblright{}.

\begin{quote}
Pendant le vol, nous avons survolé les montagnes.
\end{quote}

\begin{quote}
\emph{During the flight we flew over the mountains.}
\end{quote}

Figurative: \emph{survoler un texte} — to skim a text.

\begin{cousinweb}
\emph{survoler} is built from \emph{sur} 'over' and \emph{voler} 'to fly', which is inherited from Latin \emph{volāre} (Lewis \& Short). EtymDB does not list the compound.
\end{cousinweb}

\begin{grammarlens}[title={What you've built}]
One more word for telling, explaining and describing at B1.
\end{grammarlens}

\subsection*{Your turn}
\begin{itemize}
\raggedright
  \item \emph{Say it:} \emph{survoler}
  \item \emph{Say it:} \emph{survoler}, once more
  \item \emph{Say it:} a sentence of your own with \emph{survoler}
  \item \emph{Recall:} say the French for the landing, then the French for to land, then say \emph{survoler} again
\end{itemize}

\begin{tcolorbox}[breakable,colback=teal!4,colframe=teal!35!black,title={Before you move on}]
What does survoler mean? (\textquotedblleft{}To fly over\textquotedblright{}.) Say it once more.
\end{tcolorbox}
